\documentclass[11pt]{article}
\usepackage{amsmath, amssymb, amsthm, enumerate}
\usepackage[margin=1in, top=1.2in, bottom=1.2in]{geometry}
\usepackage{parskip}
\usepackage[colorlinks=true, linkcolor=blue, citecolor=blue, urlcolor=blue]{hyperref}
\usepackage{cleveref}
\theoremstyle{definition}
\newtheorem{definition}{Definition}[section]
\newtheorem{theorem}{Theorem}[section]
\newtheorem{lemma}{Lemma}[section]
\newtheorem{proposition}{Proposition}[section]
\newtheorem{corollary}{Corollary}[section]
\theoremstyle{remark}
\newtheorem{remark}{Remark}[section]
\newtheorem{example}{Example}[section]
\setlength{\parindent}{0pt}
\setlength{\parskip}{6pt}
\begin{document}

\begin{center}
{\Large\textbf{1.3 Subgroup, Quotient Group}}\\
\vspace{0.3cm}
{\normalsize\today}
\end{center}
\vspace{0.5cm}


% --- Page 1 ---
\begin{definition}\label{def:subgroup_topological_group}
Let $G$ be a topological group and $H \subseteq G$. Then $H$ is a \underline{subgroup of topological group $G$} if $H$ is a subgroup of group $G$ and $H$ is a closed set in topological space $G$.
\end{definition}

\begin{remark}\label{rem:induced_topology_subgroup}
$H$ is a topological group equipped with induced topology on $G$. A subgroup of an abstract group which is a topological group is a topological group. This works as restriction of continuous maps is continuous.
\end{remark}

\begin{definition}\label{def:normal_subgroup_topological_group}
A normal subgroup $N$ of topological group $G$ if $N$ is a normal subgroup of abstract group $G$.
\end{definition}

% --- Page 2 ---
\begin{proposition}\label{prop:closure_subgroup}
Let $G$ be a topological group. Let $H$ be a subgroup of abstract group $G$.

\begin{enumerate}[(i)]
\item $\overline{H}$ is a subgroup of top. group $G$.
\item If $H$ is normal subgroup of $G$ then $\overline{H}$ is normal.
\end{enumerate}
\end{proposition}

\begin{proof}
(i) Observe $\overline{H} = \bigcap_{V \in \mathcal{V}} HV$ where $\mathcal{V}$ is collection of all nbds of $e$. As $H \subseteq \overline{H}$ hence $\overline{H} \neq \emptyset$.

For $x, y \in \overline{H}$ we want to show $xy^{-1} \in \overline{H}$. We know $m: G \times G \to G$ is continuous, where $m(x,y) = xy^{-1}$. Observe $m^{-1}(\overline{H})$ is closed, and $H \times H \subseteq m^{-1}(\overline{H})$. Hence $\overline{H} \times \overline{H} = \overline{H \times H} \subseteq m^{-1}(\overline{H})$ and this implies $m(\overline{H} \times \overline{H}) \subseteq \overline{H}$.
\end{proof}

% --- Page 3 ---
---

We construct a mapping
\[
n \colon G \times G \to G, \quad (x, y) \mapsto x y x^{-1}
\]
and observe that $n$ is a continuous map. This implies that $n^{-1}(\overline{H})$ is closed and similar to (Ci).
\[
H \times H \subseteq n^{-1}(\overline{H}) \implies \overline{H} \times \overline{H} = \overline{H \times H}
\]
is a subset of $n^{-1}(\overline{H})$. Hence
\[
n(\overline{H} \times \overline{H}) \subseteq \overline{H}.
\]

\begin{definition}\label{def:left_cosets}
* Set of all left cosets set $G/H$. We have a natural map or projection
\[
\pi \colon G \to G/H, \quad x \mapsto xH.
\]
\end{definition}

\begin{remark}\label{rem:topology_on_quotient}
If $G$ is a topological group, we will topologize $G/H$ assuming $H$ is closed. Construct
\[
\mathcal{T}_{G/H} = \{ O \subseteq G/H \mid \pi^{-1}(O) \in \mathcal{T}_G \},
\]
i.e., $O$ is open if and only if $\pi^{-1}(O)$ is open in $G$.
\end{remark}

% --- Page 4 ---
---

which ensures $\pi$ is a continuous map, moreover a quotient map where $\pi$ is surjective by construction.

\begin{lemma}\label{lem:quotient_space}
$G/H$ is a $T_1$-space, and $\pi$ is open.
\end{lemma}

\begin{proof}
Since $H$ is closed, then $\alpha H$ is closed by homeomorphism. Observe $G - \alpha H$ is open in $G$.
Let $\tilde{\alpha} = \alpha H \in G/H$. Observe
$\pi^{-1}(G/H - \tilde{\alpha}) = G - \alpha H$.
Hence $\tilde{\alpha}$ is closed, $G/H$ is $T_1$.

We want to show $\pi$ is an open map.
Let $O$ be an open map in $G/H$ then
$\pi(O)$ is open in $G/H$ iff $\pi^{-1}(\pi(O))$ is open in $G$.
Now $OH = \pi^{-1}(\pi(O))$ and
$OH = \bigcup_{\alpha \in O} \alpha H$.
Hence $OH$ is open.
\end{proof}

% --- Page 5 ---
---

$\pi^{-1}(\pi(O))$ is open which implies $\pi(O)$ is open.

\begin{lemma}\label{lem:G/H_T2}
$G/H$ is a $T_2$ space.
\end{lemma}

\begin{proof}
Consider $xH \neq yH$ in $G/H$.
Now $x \in xH$ and $x \notin G \setminus yH$ which is open as $H$ is closed and $yH$ is closed.
Now, let $V$ be a neighborhood of $x$ such that $x \in V \subset G \setminus yH$.
Observe $Vx \cap yH = \emptyset$ implies $VxH \cap yH = \emptyset$ otherwise $vh_1 = gh_2$ and $v = gh_2h_1^{-1}$ implying $Vx \cap yH \neq \emptyset$.
Now, we can find a neighborhood $V_1$ of $V$ such that $V_1^{-1}V_1 \subset V$ therefore $V_1^{-1}V_1xH \cap yH = \emptyset$.
This implies $V_1xH \cap V_1yH = \emptyset$, which is
$\pi(V_1x) \cap \pi(V_1y) = \emptyset$ and by previous lemma $\pi(V_1x)$ and $\pi(V_1y)$ are open sets and $xH \in \pi(V_1x)$ and $yH \in \pi(V_1y)$.
\end{proof}

% --- Page 6 ---
\begin{remark}
We call $G/H$ the quotient space where $H$ is a closed subgroup of $H$, has equipped with quotient topology.
\end{remark}

\begin{lemma}\label{lem:quotient_topology_group}
If $G$ is a topological group, $H$ is normal subgroup, the quotient topology equipped on the quotient group $G/H$ is a topological group.
\end{lemma}

\begin{proof}
We want to show $\varphi_1: G/H \times G/H \to G/H$ defined by
\[
\varphi_1(\alpha_1, y_1) \mapsto \alpha_1 y_1^{-1}
\]
is continuous.

And $\varphi: G \times G \to G$ defined by
\[
\varphi(\alpha, y) \mapsto \alpha y^{-1}
\]
is continuous.
\end{proof}

% --- Page 7 ---
\[
\begin{array}{ccc}
G \times G & \xrightarrow{\pi \times \pi} & G/H \times G/H \\
\downarrow{\pi} & & \downarrow{\psi} \\
G & \xrightarrow{\varphi} & G/H
\end{array}
\]

Observe that $\psi_1(\pi \times \pi) = \psi \pi$ and $\pi$, $\psi$ are continuous. Hence $\psi_1(\pi \times \pi)$ are continuous.

Let $O_1$ be an open subset of $G/H$ then
$\psi_1(\pi \times \pi)^{-1}(O_1)$ is open. As $\pi \times \pi$ is an open map, observe that $\pi \times \pi (\pi \times \pi)^{-1} \psi_1^{-1}(O_1)$ is open. Hence $\psi_1^{-1}(O_1)$ is open.

\begin{proof}
$\square$
\end{proof}

\begin{lemma}
Let $G$ be a topological group, and $H$ be a subgroup.
\begin{enumerate}
    \item[(a)] If $G$ is compact then $H$ and $G/H$ are both compact.
    \item[(b)] If $G$ is locally compact then $H$ and $G/H$ have both locally compact.
\end{enumerate}
\end{lemma}

% --- Page 8 ---
\begin{proof}
\begin{enumerate}
    \item[(a)] $H$ is a closed subset of a compact set and $G/H = \pi(G)$ where $G$ is compact and $\pi$ is a continuous mapping. Hence a continuous map of a compact set is compact.

    \item[(b)] A closed subset of a locally compact set is locally compact.

    Let $S$ be locally compact, $T = \overline{S}$ and $T \subseteq S$. For $p \in T \cap S$, there exists a neighborhood of $p$ in $S$, $\overline{U_p}$ such that $\overline{U_p}$ is compact. Observe for $p \in T$, there exists an open set $T \cap U_p$ in $T$ such that $\overline{T \cap U_p} = \overline{T} \cap \overline{U_p}$ is compact as $\overline{T \cap U_p}$ is a closed subset of $\overline{U_p}$ which is compact. Hence $H$ is locally compact.

    Goal: $G/H$ is locally compact.
\end{enumerate}
\end{proof}

% --- Page 9 ---
Let $q_i \in G \setminus H$ and let $U_i$ be a neighborhood of $q_i$. Let
\[ U = \pi^{-1}(U_i), \]
and let $x \in G$ such that $\pi(x) = q_i$. As $G$ is locally compact, we get a neighborhood $O$ of $x$ such that $\overline{O}$ is compact and $\overline{O} \subseteq U$ by the proposition given below. Then we have $\pi(\overline{O}) \subseteq \pi(U) = U_i$. If $O_i = \pi(O)$ is a neighborhood of $q_i$. And $\pi(\overline{O})$ is a continuous image of a compact set, hence we get $\pi(\overline{O})$ is compact. Since $G/H$ is $T_2$-space $\pi(\overline{O})$ is closed which implies
\[ O \subseteq \overline{O} \Rightarrow O_i \subseteq \pi(\overline{O}) \quad \text{and} \quad \overline{O_i} \subseteq \pi(\overline{O}) = \pi(\overline{O}) \quad \text{which is compact, so is} \quad \overline{O_i}. \]
Thus $O_i$ is a neighborhood of $q_i$ such that $\overline{O_i}$ is compact. It follows $G/H$ is locally compact.

\begin{corollary}\label{cor:quotient_locally_compact}
If $G$ is locally compact, then $G/H$ is normal or $T_4$.
\end{corollary}

% --- Page 10 ---
\begin{proposition}\label{prop:nbd_subgroup}
For any topological group $G$, let $U$ be a neighborhood of $e$, there exists a neighborhood of $e$ $V$ such that $V \subseteq U$.
\end{proposition}

\begin{proof}
By properties (5) and (6) we know that for every neighborhood of $e$, $U$, we have a symmetric neighborhood of $e$ $V$ such that $V^2 \subseteq U$. Now $x \in \overline{V} = \bigcap_{w \in V} wV$ where $V^2$ is the set of all neighborhoods of $e$, this implies $x \in \overline{V} \subseteq V^2 \subseteq U$.
\end{proof}

\end{document}